% !TEX root = ../main.tex
% File: part1/chapters_pretraining/intro.tex
% Chương mới (cập nhật 2026) -- trục Pretraining + Optimization của LLM Research Reading List

\chapter{TIỀN HUẤN LUYỆN LLM: QUY LUẬT MỞ RỘNG, DỮ LIỆU VÀ TỐI ƯU HÓA}
\label{chap:pretraining}

Chương~\ref{chap:transformer_llms} đã cho chúng ta bản thiết kế của một LLM. Giữa một kiến trúc Transformer trên giấy và một mô hình như Llama 3 hay DeepSeek-V3 là hàng tháng trời chạy trên hàng chục nghìn GPU, tiêu tốn hàng triệu đô-la. Với chi phí như vậy, không ai có thể “thử và sai” ở quy mô thật: mọi quyết định phải được đưa ra \emph{trước} khi bấm nút huấn luyện.

Chương này bàn ba quyết định mà mọi đội tiền huấn luyện đều phải đưa ra:
\begin{enumerate}
	\item \textbf{Quy mô mô hình và lượng token:} Các \textbf{quy luật mở rộng (scaling laws)} cho phép dự đoán loss của mô hình lớn từ những thí nghiệm nhỏ (mục~\ref{sec:scaling_laws}).
	\item \textbf{Thành phần dữ liệu:} Từ The Pile đến FineWeb và DCLM, dữ liệu tiền huấn luyện chuyển dần từ “càng nhiều càng tốt” sang “càng sạch, càng có chọn lọc càng tốt” (mục~\ref{sec:pretraining_data}).
	\item \textbf{Thuật toán tối ưu:} Adam và AdamW đã thống trị gần một thập kỷ; Sophia, Adam-mini và đặc biệt là Muon đang thách thức vị thế đó (mục~\ref{sec:llm_optimizers}).
\end{enumerate}
Cuối chương, chúng ta ghép tất cả lại qua hai “công thức” huấn luyện được công bố chi tiết nhất: Llama 3 và DeepSeek-V3 (mục~\ref{sec:pretraining_case_studies}).
